%%% Основные сведения %%%
\newcommand{\thesisAuthorLastName}{Леменкова}
\newcommand{\thesisAuthorOtherNames}{Полина Алексеевна}
\newcommand{\thesisAuthorInitials}{П.\,А.}
\newcommand{\thesisAuthor}           % Диссертация, ФИО автора
{%
    \texorpdfstring{% \texorpdfstring takes two arguments and uses the first for (La)TeX and the second for pdf
        \thesisAuthorLastName~\thesisAuthorOtherNames% так будет отображаться на титульном листе или в тексте, где будет использоваться переменная
    }{%
        \thesisAuthorLastName, \thesisAuthorOtherNames% эта запись для свойств pdf-файла. В таком виде, если pdf будет обработан программами для сбора библиографических сведений, будет правильно представлена фамилия.
    }
}
\newcommand{\thesisAuthorShort}        % Диссертация, ФИО автора инициалами
{\thesisAuthorInitials~\thesisAuthorLastName}
%\newcommand{\thesisUdk}                % Диссертация, УДК
%{\fixme{xxx.xxx}}
\newcommand{\thesisTitle}              % Диссертация, название
{Морфоструктурные особенности глубоководных желобов Тихого океана, проблема их происхождения}
\newcommand{\thesisSpecialtyNumber}    % Диссертация, специальность, номер
{25.00.03}
\newcommand{\thesisSpecialtyTitle}     % Диссертация, специальность, название (название взято с сайта ВАК для примера)
{Геотектоника и геодинамика}

\newcommand{\thesisDegree}             % Диссертация, ученая степень
{кандидата геолого-минералогических наук}
\newcommand{\thesisDegreeShort}        % Диссертация, ученая степень, краткая запись
{канд. геол.-мин. наук}
\newcommand{\thesisCity}               % Диссертация, город написания диссертации
{Москва}
\newcommand{\thesisYear}               % Диссертация, год написания диссертации
{\the\year}
\newcommand{\thesisOrganization}       % Диссертация, организация
{\fixme{Федеральное государственное автономное образовательное учреждение высшего
образования <<Длинное название образовательного учреждения <<АББРЕВИАТУРА>>}}
\newcommand{\thesisOrganizationShort}  % Диссертация, краткое название организации для доклада
{\fixme{НазУчДисРаб}}

\newcommand{\thesisInOrganization}     % Диссертация, организация в предложном падеже: Работа выполнена в ...
{Лаборатории № 303 Региональной геофизики и природных катастроф Отделении III: Природных, природно-техногенных катастроф и сейсмичности Земли Института физики Земли им. О. Шмидта Российской Академии Наук (ИФЗ РАН)}

%% \newcommand{\supervisorDead}{}           % Рисовать рамку вокруг фамилии
\newcommand{\supervisorFio}              % Научный руководитель, ФИО
{Илларионов Владимир Константинович}
\newcommand{\supervisorRegalia}          % Научный руководитель, регалии
{кандидат геолого-минералогических наук, ведущий научный сотрудник Лаборатории № 303 Региональной геофизики и природных катастроф ИФЗ им. О. Шмидта РАН}
\newcommand{\supervisorFioShort}         % Научный руководитель, ФИО
{\fixme{И.\,О.~Фамилия}}
\newcommand{\supervisorRegaliaShort}     % Научный руководитель, регалии
{\fixme{уч.~ст.,~уч.~зв.}}

%% \newcommand{\supervisorTwoDead}{}        % Рисовать рамку вокруг фамилии
%% \newcommand{\supervisorTwoFio}           % Второй научный руководитель, ФИО
%% {\fixme{Фамилия Имя Отчество}}
%% \newcommand{\supervisorTwoRegalia}       % Второй научный руководитель, регалии
%% {\fixme{уч. степень, уч. звание}}
%% \newcommand{\supervisorTwoFioShort}      % Второй научный руководитель, ФИО
%% {\fixme{И.\,О.~Фамилия}}
%% \newcommand{\supervisorTwoRegaliaShort}  % Второй научный руководитель, регалии
%% {\fixme{уч.~ст.,~уч.~зв.}}

\newcommand{\opponentOneFio}           % Оппонент 1, ФИО
{Хортов Алексей Владимирович}
\newcommand{\opponentOneRegalia}       % Оппонент 1, регалии
{доктор геолого-минералогических наук}
\newcommand{\opponentOneJobPlace}      % Оппонент 1, место работы
{Институт океанологии им. П.П.  Ширшова РАН, Лаборатория геодинамики, георесурсов и геоэкологии ФГБУН}
\newcommand{\opponentOneJobPost}       % Оппонент 1, должность
{ведущий научный сотрудник}

\newcommand{\opponentTwoFio}           % Оппонент 2, ФИО
{Пискарев-Васильев Алексей Лазаревич}
\newcommand{\opponentTwoRegalia}       % Оппонент 2, регалии
{доктор геолого-минералогических наук}
\newcommand{\opponentTwoJobPlace}      % Оппонент 2, место работы
{Институт Наук о Земле СПГУ, кафедра геофизики, ФГБУ ВНИИОкеангеология им. акад. И.С. Грамберга}
\newcommand{\opponentTwoJobPost}       % Оппонент 2, должность
{профессор, академик РАЕН}

%% \newcommand{\opponentThreeFio}         % Оппонент 3, ФИО
%% {\fixme{Фамилия Имя Отчество}}
%% \newcommand{\opponentThreeRegalia}     % Оппонент 3, регалии
%% {\fixme{кандидат физико-математических наук}}
%% \newcommand{\opponentThreeJobPlace}    % Оппонент 3, место работы
%% {\fixme{Основное место работы c длинным длинным длинным длинным названием}}
%% \newcommand{\opponentThreeJobPost}     % Оппонент 3, должность
%% {\fixme{старший научный сотрудник}}

\newcommand{\leadingOrganizationTitle} % Ведущая организация, дополнительные строки. Удалить, чтобы не отображать в автореферате
{Геологический Институт РАН (ГИН РАН)}

\newcommand{\defenseDate}              % Защита, дата
{DD октября 2021~г.~в~XX часов}
\newcommand{\defenseCouncilNumber}     % Защита, номер диссертационного совета
{Д\,002.001.01}
\newcommand{\defenseCouncilTitle}      % Защита, учреждение диссертационного совета
{Институт физики Земли им. О. Шмидта Российской Академии Наук}
\newcommand{\defenseCouncilAddress}    % Защита, адрес учреждение диссертационного совета
{123242, г. Москва, Б. Грузинская ул., д. 10, стр. 1}
\newcommand{\defenseCouncilPhone}      % Телефон для справок
{+7~(0000)~00-00-00}

\newcommand{\defenseSecretaryFio}      % Секретарь диссертационного совета, ФИО
{Камзолкин Владимир Анатольевич}
\newcommand{\defenseSecretaryRegalia}  % Секретарь диссертационного совета, регалии
{к. г.-м. н.}         % Для сокращений есть ГОСТы, например: ГОСТ Р 7.0.12-2011 + http://base.garant.ru/179724/#block_30000

\newcommand{\synopsisLibrary}          % Автореферат, название библиотеки
{ИФЗ РАН}
\newcommand{\synopsisDate}             % Автореферат, дата рассылки
{\fixme{DD} июня \the\year~года}

% To avoid conflict with beamer class use \providecommand
\providecommand{\keywords}%            % Ключевые слова для метаданных PDF диссертации и автореферата
{}
